\documentclass{article}
\usepackage[utf8]{inputenc}

\usepackage{amsmath}
\usepackage{amssymb}
\usepackage{enumitem}
\usepackage{marvosym}
\usepackage{amscd}

\usepackage{csquotes}

\title{Strict Order}
\author{Adam W. Grant}
\date{January 2021}

\begin{document}

\maketitle

The textbook \cite{davey_priestley_2002} explains a strict order as follows
\begin{displayquote}
An order relation $\leq$ on $P$ gives rise to a relation $<$ of \textbf{strict inequality}: $x < y$ if and only if $x \leq y$ and $x\neq y$.
\end{displayquote}
The following calculation shows how to say the same without points
\begin{itemize}
    \item $x < y$
    \item[$\equiv$] $\{\text{Definition of $<$}\}$
    \item[] $x \leq y \wedge x \neq y$
    \item[$\equiv$] $\{\text{Expand abbreviation}\}$
    \item[] $x \leq y \wedge \neg (x = y)$
    \item[$\equiv$] $\{\text{Identity}\}$
    \item[] $x \leq y \wedge \neg (x \,id\, y)$
    \item[$\equiv$] $\{\text{Difference}\}$
    \item[] $x (\leq - \;id) y$
    \item[\Heart]
\end{itemize}
So the pointfree equivalent of the above definition is
\begin{equation*}
    < \;= \; \leq - \;id
\end{equation*}
Informally, we can say that removing all of the reflexive arrows from a partial order $\leq$ gives us its corresponding strict order $<$.

\bibliographystyle{plain}
\bibliography{refs}

\end{document}
